\documentclass[11pt]{article}

\usepackage[utf8]{inputenc}
\usepackage[T1]{fontenc}
\usepackage{lmodern}
\usepackage{amsmath,amssymb,amsthm,mathtools}
\usepackage{booktabs}
\usepackage{enumitem}
\usepackage[margin=1in]{geometry}
\usepackage{microtype}
\usepackage{xcolor}
\usepackage[colorlinks=true,linkcolor=blue!55!black,citecolor=blue!55!black,urlcolor=blue!70!black]{hyperref}

\newtheorem{theorem}{Theorem}[section]
\newtheorem{lemma}[theorem]{Lemma}
\newtheorem{proposition}[theorem]{Proposition}
\newtheorem{corollary}[theorem]{Corollary}
\theoremstyle{definition}
\newtheorem{definition}[theorem]{Definition}
\newtheorem{problem}[theorem]{Problem}
\theoremstyle{remark}
\newtheorem{remark}[theorem]{Remark}

\newcommand{\R}{\mathbb{R}}
\newcommand{\N}{\mathbb{N}}
\newcommand{\Deltaop}{\Delta}
\newcommand{\ind}{\mathbf{1}}

\title{\textbf{Extremal Multiplicity Profiles for Integer Horizontal Chords}\\
\large Classification and Polynomial Realization}

\author{Zixuan Ye \and GPT-6 Astra}

\date{September 2026}

\begin{document}

\maketitle

\begin{abstract}
Let \(F:\R\to\R\) be continuous.  For \(h>0\), write
\[
\Delta_hF(x)=F(x+h)-F(x),\qquad
\nu_h(F)=\#\{x\in\R:\Delta_hF(x)=0\},
\]
where zeros are counted as distinct real points.  For an integer \(n\ge1\), put
\[
m_n(F)=\nu_n(F),\quad M(F)=\sum_{n\ge1}m_n(F).\]
A 2025 problem from the final round of the All-Russian Mathematical Olympiad
asks, for a prescribed chord of length \(2025\), for the least possible finite
value of \(M(F)\); the official answer is \(4049\).

We develop the equality theory for an arbitrary target \(N\).
First, \(m_N(F)>0\) and \(M(F)<\infty\) imply
\[
M(F)\ge 2N-1.
\]
We then completely classify the multiplicity vectors at equality.  They are
exactly the two families
\[
V_{N,r}=(4^r,2^{\,N-2r-1},0^r,1),
\qquad 0\le r<N/2,
\]
and
\[
W_{N,r}=(4^r,3,2^{\,N-2r-3},1,0^r,1),
\qquad r\ge0,\quad 2r+3\le N.
\]
Here \(a^k\) denotes \(k\) consecutive entries equal to \(a\).
Consequently the number of extremal vectors is
\[
2\left\lceil\frac N2\right\rceil-1,
\]
equal to \(N\) for odd \(N\) and \(N-1\) for even \(N\).
The transverse extremal vectors are precisely the \(V_{N,r}\).

The realization problem is also solved completely.  We construct a
two-parameter \(C^2\) model by prescribing three inverse branches and their
horizontal separations, obtaining the exact real-length formula
\[
\nu_h
=
2\ind_{\{h<c\}}+2\ind_{\{h<e\}}
+\ind_{\{h=c\}}+\ind_{\{h=e\}}+\ind_{\{h=N\}},
\qquad c+e=N.
\]
A positivity-preserving polynomialization, based on a fixed critical-point
factor and two exact moment constraints, then shows that every \(V_{N,r}\)
and every \(W_{N,r}\) is realizable by a polynomial with rational
coefficients.  Thus the extremal support, the nontransverse multiplicity
allocation, and polynomial existence are all determined for every target.
\end{abstract}

\tableofcontents

\section{Introduction}

A horizontal chord of a real-valued function is a horizontal line segment
whose endpoints lie on the graph.  Equivalently, a chord of length \(h>0\)
is a solution of
\[
F(x+h)=F(x).
\]
The classical horizontal-chord theory asks which \emph{lengths} occur.
The universal chord theorem of L\'evy and the later theorem of Hopf are
fundamental results in this direction; modern accounts include
Oxtoby~\cite{Oxtoby1972}, Burns--Davidovich--Davis~\cite{BurnsDavidovichDavis2017},
and Davis--Troubetzkoy~\cite{DavisTroubetzkoy2024}.  Recent work of
Ciudin--Ionascu~\cite{CiudinIonascu2026} continues the study of the topology
and isolated points of chord sets.

The present problem is different.  We restrict attention to \emph{integer}
lengths and ask not only whether a length occurs, but how many distinct
chords of that length occur.  Thus the basic object is the multiplicity
sequence
\[
(m_1(F),m_2(F),\ldots).
\]
This refinement was prompted directly by Problem 11.8 of the 2025 final
round of the All-Russian Mathematical Olympiad, proposed by A.\,S.~Kuznetsov.
For a continuous \(F:\R\to\R\) with finitely many integer horizontal chords
and with at least one chord of length \(2025\), the problem asks for the
smallest possible total number of chords.  The official solution proves
that the answer is \(4049\)~\cite{VSO2025}.

The first part of the olympiad argument extends naturally from \(2025\) to
an arbitrary target \(N\), giving the sharp minimum \(2N-1\).  The purpose
of this paper is to understand the much more rigid equality case.  Three
questions drive the analysis:

\begin{enumerate}[label=(\roman*)]
\item Which integer lengths may be absent when \(M(F)=2N-1\)?
\item Once the support is known, which multiplicities can actually occur
without a transversality assumption?
\item Can every admissible extremal vector be realized by a polynomial,
uniformly in \(N\)?
\end{enumerate}

All three questions admit complete answers.  The missing lengths below
\(N\) form a terminal interval.  Global compatibility among the difference
functions then reduces the apparent local choices \((1,3),(2,2),(3,1)\)
to a single possible exceptional pair, yielding the two families \(V\) and
\(W\).  Finally, an inverse-branch construction realizes all of them at
once, and a positive polynomial weight with two exact moments converts the
construction into one with rational polynomial coefficients.

\subsection{Relation to the existing literature}

The distinction between chord \emph{support} and chord \emph{multiplicity}
is important.  Hopf-type results characterize, or constrain, the set
\[
\mathcal H(F)=\{h>0:\nu_h(F)>0\},
\]
whereas the present work studies the cardinalities of the fibers
\(\{x:F(x+h)=F(x)\}\) at integer \(h\), under a finite-total constraint.
The papers cited above are therefore closely related background, but their
principal object is different.

The closest direct predecessor we located is the 2025 olympiad problem and
its official solution, which proves the sharp total for the single target
\(2025\) and contains the complementary-length mechanism used below.
A literature search conducted through September 2026 located no published
paper or arXiv preprint whose main object is the finite integer multiplicity
sequence \(n\mapsto m_n(F)\), nor one containing the extremal \(V/W\)
classification or the all-parameter polynomial realization theorem proved
here.  This is a report of the search outcome, not a claim of absolute
priority over all possible literature.

\section{Definitions and elementary identities}

\begin{definition}
For a continuous function \(F:\R\to\R\) and \(h>0\), define
\[
\Delta_hF(x):=F(x+h)-F(x),
\qquad
\nu_h(F):=\#\{x\in\R:\Delta_hF(x)=0\}.
\]
If the zero set is infinite, we write \(\nu_h(F)=\infty\).
For \(n\in\N\),
\[
m_n(F):=\nu_n(F),
\qquad
M(F):=\sum_{n\ge1}m_n(F)\in\N\cup\{\infty\}.
\]
\end{definition}

\begin{definition}
Fix \(N\in\N\).  We call \(F\) \emph{\(N\)-extremal} if
\[
m_N(F)>0,\qquad M(F)=2N-1.
\]
If \(F\in C^1(\R)\), we call it \emph{transverse below \(N\)} if every
zero of \(\Delta_nF\) is simple for \(1\le n<N\).
\end{definition}

We repeatedly use the cocycle identities
\begin{align}
\Delta_{a+b}F(x)
&=\Delta_aF(x)+\Delta_bF(x+a), \label{eq:cocycle}\\
\Delta_nF(x)
&=\sum_{j=0}^{n-1}\Delta_1F(x+j). \label{eq:unit-sum}
\end{align}

\begin{lemma}[Common tail orientation]\label{lem:tails}
If \(M(F)<\infty\), then after replacing \(F\) by \(-F\), if necessary,
\[
\Delta_1F(x)>0
\]
for all sufficiently large \(|x|\).  Consequently, for every fixed
\(n\ge1\), \(\Delta_nF(x)>0\) on both sufficiently remote tails.
\end{lemma}

\begin{proof}
Since \(m_1(F)<\infty\), the continuous function
\(g=\Delta_1F\) is nonzero and has a constant sign on each sufficiently
remote tail.  If those two signs were opposite, then
\eqref{eq:unit-sum} would give opposite tail signs for every
\(\Delta_nF\).  The intermediate value theorem would then imply
\(m_n(F)>0\) for every \(n\), contradicting \(M(F)<\infty\).
Thus the two tail signs agree, and a global sign change of \(F\) makes
them positive.  Equation~\eqref{eq:unit-sum} gives the assertion for
each fixed \(n\).
\end{proof}

Two immediate consequences will be used without further comment:
\[
m_d(F)=0\Longrightarrow \Delta_dF>0,
\qquad
m_d(F)=1\Longrightarrow \Delta_dF\ge0,
\tag{2.1}
\]
after the orientation of Lemma~\ref{lem:tails}.

\section{The sharp total and equality constraints}

\begin{theorem}[Complementary-length inequality]\label{thm:pair}
Suppose \(M(F)<\infty\), \(m_N(F)>0\), and \(a,b\in\N\) satisfy
\(a+b=N\).  Then
\[
m_a(F)+m_b(F)\ge4.
\]
\end{theorem}

\begin{proof}
Translate the variable and subtract a constant so that one target chord is
\[
F(0)=F(N)=0.
\]
Orient the tails by Lemma~\ref{lem:tails}.  Suppose, to the contrary, that
\(m_a+m_b\le3\).  Interchanging \(a,b\), we may assume \(m_a\le1\).
By \eqref{eq:unit-sum} and the positive tails, a negative value of
\(\Delta_aF\) would force at least two zeros; hence
\[
\Delta_aF\ge0.
\tag{3.1}
\]

If \(a=b\), then
\[
0=\Delta_NF(0)=\Delta_aF(0)+\Delta_aF(a),
\]
so both terms vanish, contradicting \(m_a\le1\).  Hence \(a\ne b\).

Let \(q=N/\gcd(a,N)\).  For \(1\le j\le q-1\), let \(r_j\) be the
representative of \(ja\bmod N\) in \(\{1,\ldots,N-1\}\).  Choose
\[
r_0=r_q=
\begin{cases}
0,&\Delta_aF(0)>0,\\
N,&\Delta_aF(0)=0.
\end{cases}
\]
Every step \(r_j\to r_{j+1}\) either increases by \(a\) or decreases by
\(b=N-a\).  Put
\[
d_j:=F(r_{j+1})-F(r_j).
\]
An increasing step gives \(d_j=\Delta_aF(r_j)\ge0\).  At least one such
step is strictly positive: it is the first step when \(r_0=0\), and the
last step when \(r_0=N\), because in the latter case the unique possible
zero of \(\Delta_aF\) is already at \(0\).

Since
\[
\sum_{j=0}^{q-1}d_j=F(r_q)-F(r_0)=0,
\]
some decreasing step has \(d_j<0\).  If its landing point is \(t\), then
\[
\Delta_bF(t)>0,
\qquad 0<t<a.
\tag{3.2}\label{eq:3.2}
\]
On the other hand,
\[
\Delta_bF(0)=-\Delta_aF(b)\le0,
\qquad
\Delta_bF(a)=-\Delta_aF(0)\le0.
\tag{3.3}\label{eq:3.3}
\]
If both values in \eqref{eq:3.3} are negative, then the positive tails,
the two negative endpoint values, and \eqref{eq:3.2} force four zeros of
\(\Delta_bF\).  If exactly one endpoint value is zero, then
\(\Delta_bF\) has at least three zeros and \(\Delta_aF\) has at least one.
Both endpoint values cannot vanish, for that would give two distinct
zeros of \(\Delta_aF\).  In all cases \(m_a+m_b\ge4\), a contradiction.
\end{proof}

\begin{theorem}[Sharp total]\label{thm:sharp}
If \(m_N(F)>0\) and \(M(F)<\infty\), then
\[
M(F)\ge 2N-1.
\]
The bound is sharp.
\end{theorem}

\begin{proof}
Pair the lengths \(1,\ldots,N-1\) by \(a\leftrightarrow N-a\) and apply
Theorem~\ref{thm:pair}.  For odd \(N\), the \((N-1)/2\) pairs contribute
at least \(4\) each, and the target contributes at least \(1\).  For even
\(N\), the unequal pairs contribute \(4\) each, while \(a=N/2\) contributes
at least \(2\).  In both cases the total is at least \(2N-1\).

For sharpness, let
\[
P_N(x)=x^3-\frac{N^2}{4}x.
\]
Then
\[
\frac{\Delta_nP_N(x)}{n}
=
3\left(x+\frac n2\right)^2+\frac{n^2-N^2}{4}.
\]
Thus \(m_n(P_N)=2\) for \(1\le n<N\), \(m_N(P_N)=1\), and
\(m_n(P_N)=0\) for \(n>N\).  Hence \(M(P_N)=2N-1\).
\end{proof}

Therefore every \(N\)-extremal function satisfies
\begin{equation}
m_N=1,\qquad
m_n=0\ \ (n>N),\qquad
m_a+m_{N-a}=4\ \ (1\le a<N).
\tag{E}\label{eq:E}
\end{equation}

\section{Missing lengths and terminal support}

The missing lengths have an additive structure even before extremality.

\begin{theorem}[Missing-length semigroup]\label{thm:semigroup}
If \(M(F)<\infty\), then
\[
\Gamma(F):=\{0\}\cup\{n\in\N:m_n(F)=0\}
\]
is a numerical semigroup.
\end{theorem}

\begin{proof}
Only additive closure requires proof.  If \(a,b\) are missing, then
\(\Delta_aF\) and \(\Delta_bF\) never vanish and therefore each has a
constant sign.  Their signs must agree, since
\[
\Delta_{ab}F(x)
=
\sum_{j=0}^{b-1}\Delta_aF(x+ja)
=
\sum_{j=0}^{a-1}\Delta_bF(x+jb)
\]
cannot simultaneously have opposite signs.  Hence
\[
\Delta_{a+b}F(x)=\Delta_aF(x)+\Delta_bF(x+a)
\]
never vanishes, so \(a+b\) is missing.  Since \(M(F)<\infty\), only
finitely many positive integers occur, and \(\Gamma(F)\) is cofinite.
\end{proof}

For the equality case, much more is true.

\begin{lemma}[Strict block propagation]\label{lem:strict-block}
Assume \(F(0)=F(N)\), \(M(F)<\infty\), and let \(1\le d<N\) be missing.
Write
\[
N=qd+s,\qquad 0\le s<d.
\]
Then \(s\ne0\).  If \(q\ge2\), then
\[
m_1(F)\ge2(q+1).
\]
\end{lemma}

\begin{proof}
Orient the tails positively and set \(g=\Delta_1F\),
\(g_j=g(j)\).  Since \(\Delta_dF>0\), if \(s=0\) then
\(\Delta_NF(0)\) is a sum of \(q\) positive \(d\)-blocks, impossible.
Thus \(1\le s<d\).

For \(t=0,\ldots,q\), let
\[
R_t=\{td,\ldots,td+s-1\}.
\]
The complement of \(R_t\) in \(\{0,\ldots,N-1\}\) is a disjoint union
of \(q\) consecutive \(d\)-blocks, each of positive sum.  Since
\(\sum_{j=0}^{N-1}g_j=0\), every \(R_t\) has negative sum.
Between \(R_t\) and \(R_{t+1}\) lies
\[
G_t=\{td+s,\ldots,(t+1)d-1\},
\qquad 0\le t<q.
\]
Because \(R_t\cup G_t\) is a positive \(d\)-block, \(G_t\) has positive
sum.

Choose one negative sample from each \(R_t\) and one positive sample from
each \(G_t\).  These chosen signs alternate, beginning and ending
negative.  Together with the positive tails, the intermediate value
theorem forces \(2(q+1)\) distinct zeros of \(g\).
\end{proof}

\begin{lemma}[Weak gap filling]\label{lem:weak-gap}
Assume \(F(0)=F(N)\), \(M(F)<\infty\), and the positive tail orientation.
Let
\[
2\le k<s<N/2.
\]
Suppose
\[
m_1=4,\qquad
m_{N-j}=0\ \ (1\le j<k),\qquad
\Delta_{N-s}F\ge0.
\]
Then
\[
m_k\ge4.
\]
\end{lemma}

\begin{proof}
Put \(g=\Delta_1F\), \(g_j=g(j)\), and define
\[
P_u=\sum_{j=0}^{u-1}g_j,
\qquad
Q_u=\sum_{j=N-u}^{N-1}g_j.
\]
The target equation gives \(\sum_{j=0}^{N-1}g_j=0\).  If
\(\Delta_{N-u}F\ge0\), evaluation at \(0\) and at \(u\) gives
\[
P_u\le0,\qquad Q_u\le0,
\tag{4.1}\label{eq:PQweak}
\]
with strict inequalities when \(N-u\) is missing.  Hence
\[
P_j,Q_j<0\quad(1\le j<k),
\qquad
P_s,Q_s\le0.
\tag{4.2}
\]
In particular \(g_0,g_{N-1}<0\).

We claim \(P_k<0\).  Suppose \(P_k\ge0\).  Since \(P_{k-1}<0\),
we have \(g_{k-1}>0\).  Also
\[
\sum_{j=s}^{N-1}g_j=-P_s\ge0,
\]
and \(g_{N-1}<0\), so some \(g_t>0\) with \(t\ge s\).
The two positive samples \(g_{k-1}\) and \(g_t\), together with the
negative endpoint samples and the positive tails, already account for
the four zeros allowed by \(m_1=4\).  Therefore every sampled value
between those two positive samples must itself be strictly positive;
otherwise an additional zero or two additional sign changes would occur.
Hence \(g_k,\ldots,g_{s-1}>0\), and
\[
P_s=P_k+\sum_{j=k}^{s-1}g_j>0,
\]
contradicting \eqref{eq:PQweak}.  Thus \(P_k<0\).  Reversing the sample
order gives \(Q_k<0\).

Now set
\[
T_j=\sum_{\ell=0}^{k-1}g_{j+\ell}
=\Delta_kF(j),
\qquad 0\le j\le N-k.
\]
Then
\[
T_0=P_k<0,\qquad T_{N-k}=Q_k<0.
\]
Counting the number of times each sample occurs in the sliding sums,
\[
\sum_{j=0}^{N-k}T_j
=
-\sum_{u=1}^{k-1}(P_u+Q_u)>0.
\]
Thus some interior \(T_j\) is positive.  Since \(\Delta_kF\) has positive
tails, the sampled sign pattern negative--positive--negative forces four
distinct zeros.  Hence \(m_k\ge4\).
\end{proof}

\begin{theorem}[Terminal support]\label{thm:terminal}
Let \(F\) be \(N\)-extremal.  Then there is a unique integer
\(r\) with \(0\le r<N/2\) such that
\[
\{1\le d<N:m_d=0\}
=
\{N-r,N-r+1,\ldots,N-1\}.
\tag{4.3}\label{eq:terminal}
\]
The set on the right is empty when \(r=0\).
\end{theorem}

\begin{proof}
Let \(A=\{1\le d<N:m_d=0\}\).  If \(d<N/2\) belongs to \(A\),
Lemma~\ref{lem:strict-block} gives \(m_1\ge6\), contradicting
\eqref{eq:E}.  The value \(d=N/2\) is also impossible, because a missing
half-length would make \(\Delta_NF\) a sum of two positive half-length
differences.  Therefore \(A\subset(N/2,N)\).

If \(A=\varnothing\), take \(r=0\).  Otherwise choose \(d\in A\).
Lemma~\ref{lem:strict-block}, now with \(q=1\), gives \(m_1\ge4\);
hence \eqref{eq:E} forces
\[
m_1=4,\qquad m_{N-1}=0.
\]
Let \(\min A=N-r\).  Then \(1\le r<N/2\).
We prove inductively that \(N-k\in A\) for \(1\le k\le r\).
The case \(k=1\) is known.  If \(2\le k<r\) and
\(N-j\in A\) for every \(j<k\), apply
Lemma~\ref{lem:weak-gap} with \(s=r\).  Since \(N-r\) is missing,
\(\Delta_{N-r}F>0\), so \(m_k\ge4\).  The equality
\(m_k+m_{N-k}=4\) then gives \(m_{N-k}=0\).  Finally \(N-r\in A\)
by definition.  Minimality of \(N-r\) excludes every other missing
length below it.
\end{proof}

\section{Single zeros and the complete extremal classification}

Terminal support still leaves a formal possibility: in a central
complementary pair, equality in \eqref{eq:E} allows
\[
(1,3),\quad(2,2),\quad(3,1).
\]
The next lemma shows that a single zero is highly constrained.

\begin{lemma}[Weak block propagation for a single zero]\label{lem:single}
Assume \(F(0)=F(N)\), \(M(F)<\infty\), and \(m_d=1\) for some \(d<N\).

\begin{enumerate}[label=(\alph*)]
\item If \(d\mid N\), a contradiction follows.
\item If \(q=\lfloor N/d\rfloor\ge2\), then
\[
m_1\ge2(q+1).
\]
\item If \(d>N/2\) and \(m_1\le3\), then
\[
d=N-1,\qquad m_1=3.
\]
\end{enumerate}
\end{lemma}

\begin{proof}
Orient the tails positively.  Since \(m_d=1\),
\(\Delta_dF\ge0\).

If \(N=qd\), the target identity is a sum of \(q\ge2\) nonnegative
\(d\)-differences.  For the sum to vanish, all \(q\) terms must vanish
at distinct starting points, contradicting \(m_d=1\).  This proves (a).

For (b), write \(N=qd+s\), \(1\le s<d\), and repeat the block notation
from Lemma~\ref{lem:strict-block}.  The complement of each remainder
block \(R_t\) is a union of \(q\) nonnegative \(d\)-blocks.  At most one
of those blocks can have zero sum, while \(q\ge2\), so the complement
has strictly positive sum and \(R_t\) has strictly negative sum.
Because \(R_t\cup G_t\) is a nonnegative \(d\)-block, \(G_t\) has
strictly positive sum.  The same alternating sign argument yields
\(m_1\ge2(q+1)\).

For (c), write \(N=d+s\), \(1\le s<d\).  The two remainder blocks
\[
R_0=\{0,\ldots,s-1\},
\qquad
R_1=\{d,\ldots,N-1\}
\]
have nonpositive sums, and at least one has strictly negative sum;
otherwise \(\Delta_dF\) would vanish at two distinct starting points.
The middle block
\[
G=\{s,\ldots,d-1\}
\]
therefore has positive sum.  If both remainder blocks contain a negative
sample, then negative--positive--negative, together with the positive
tails, forces \(m_1\ge4\).  Thus, under \(m_1\le3\), one remainder block
contains no negative sample.  Since its sum is nonpositive, all its
\(s\) samples are zero.  The positive middle block and the negative
other remainder block force two further zeros of \(\Delta_1F\), distinct
from those \(s\) sampled zeros.  Hence \(m_1\ge s+2\).  Therefore
\(s=1\) and \(m_1=3\), i.e. \(d=N-1\).
\end{proof}

\begin{corollary}\label{cor:single-upper}
For an \(N\)-extremal function, a coordinate \(m_d=1\) with \(d<N\)
can occur only for
\[
N/2<d<N.
\]
\end{corollary}

\begin{proof}
If \(d\le N/2\), either \(d\mid N\) or
\(\lfloor N/d\rfloor\ge2\).  Lemma~\ref{lem:single} then contradicts
\(m_1\le4\), which follows from \eqref{eq:E}.
\end{proof}

We can now state the main classification theorem.

\begin{theorem}[Complete extremal multiplicity classification]
\label{thm:classification}
Let \(F:\R\to\R\) be continuous and \(N\)-extremal.
Then its vector
\[
(m_1(F),\ldots,m_N(F))
\]
is exactly one of the following two forms:
\[
V_{N,r}
=
(4^r,2^{\,N-2r-1},0^r,1),
\qquad 0\le r<N/2,
\tag{5.1}\label{eq:V}
\]
or
\[
W_{N,r}
=
(4^r,3,2^{\,N-2r-3},1,0^r,1),
\qquad r\ge0,\quad 2r+3\le N.
\tag{5.2}\label{eq:W}
\]
\end{theorem}

\begin{proof}
Let \(r\) be given by Theorem~\ref{thm:terminal}.  The missing terminal
block forces, by \eqref{eq:E},
\[
m_n=4\quad(1\le n\le r),
\qquad
m_n=0\quad(N-r\le n<N).
\]
Every central coordinate \(r<n<N-r\) is positive.  In a central
complementary pair, the two multiplicities are positive integers summing
to \(4\).  By Corollary~\ref{cor:single-upper}, the only possibilities,
written in short-length/long-length order, are
\[
(2,2)\quad\text{or}\quad(3,1).
\tag{5.3}
\]

First suppose \(r=0\).  If no single-zero coordinate occurs, every pair
is \((2,2)\), giving \(V_{N,0}\).  If a single-zero coordinate occurs,
Lemma~\ref{lem:single}(c) forces it to be \(N-1\) and forces
\(m_1=3\).  Hence the unique exceptional pair is \((3,1)\), giving
\(W_{N,0}\).

Now suppose \(r\ge1\), so \(m_1=4\).  Assume a central long length
\(d=N-s\) has \(m_d=1\).  Since \(d\) is central,
\[
r<s<N/2.
\]
If \(s\ge r+2\), choose \(k=r+1<s\).  The terminal support gives
\(m_{N-j}=0\) for every \(1\le j<k\), while \(m_{N-s}=1\) gives
\(\Delta_{N-s}F\ge0\).  Lemma~\ref{lem:weak-gap} yields
\(m_{r+1}\ge4\), and then \eqref{eq:E} forces
\[
m_{N-r-1}=0,
\]
contradicting the fact that the missing block starts at \(N-r\).
Therefore \(s=r+1\).  Thus at most one exceptional \((3,1)\) pair can
occur, and it must be exactly
\[
(r+1,N-r-1).
\]
The two lengths are distinct precisely when \(2r+3\le N\).  This gives
\(W_{N,r}\); otherwise all central pairs are \((2,2)\), giving
\(V_{N,r}\).
\end{proof}

\begin{corollary}[Number of extremal vectors]\label{cor:count}
The exact number of \(N\)-extremal multiplicity vectors is
\[
\left\lceil\frac N2\right\rceil
+
\left\lfloor\frac{N-1}{2}\right\rfloor
=
2\left\lceil\frac N2\right\rceil-1.
\]
Thus it equals \(N\) when \(N\) is odd and \(N-1\) when \(N\) is even.
\end{corollary}

The theorem above proves necessity.  The rest of the paper proves that
every listed vector is realizable, even polynomially.

\section{An inverse-branch model with exact real chord counts}

We first construct a \(C^2\) function whose \emph{entire real-length}
chord count is explicit.

Fix real numbers
\[
c,e\ge\frac12,\qquad c+e=N.
\]
Set
\[
\delta=\frac18,\qquad
\alpha=\frac1{16},\qquad
\gamma=\frac{\delta}{\sqrt3},
\qquad
v=c-\delta,\qquad w=c+\delta.
\]
For \(b\ge1/2\), define
\[
I_b(t)
=
\alpha\sqrt t
+\gamma\bigl(1-\sqrt{1-t}\bigr)
+\bigl(b-\delta-\alpha-\gamma\bigr)t,
\qquad 0\le t\le1,
\]
and
\[
\theta(t)=1-\sqrt{1-t^{1/3}}.
\]
The coefficient \(b-\delta-\alpha-\gamma\) is positive, hence
\(I_b\) is strictly increasing, with
\[
I_b(0)=0,\qquad I_b(1)=b-\delta.
\]

Define \(H_{N,c}:\R\to\R\) by
\[
H_{N,c}(x)=
\begin{cases}
-(x/\alpha)^2,
&x\le0,\\[1mm]
-I_c^{-1}(x),
&0\le x\le v,\\[1mm]
\operatorname{sgn}(x-c)
\left(
\dfrac{2|x-c|}{\delta}
-\left(\dfrac{x-c}{\delta}\right)^2
\right)^3,
&v\le x\le w,\\[3mm]
I_e^{-1}(N-x),
&w\le x\le N,\\[1mm]
((x-N)/\alpha)^2,
&x\ge N.
\end{cases}
\tag{6.1}\label{eq:H}
\]

\begin{proposition}[Shape and regularity]\label{prop:Hshape}
The function \(H_{N,c}\) is \(C^2\).  Its strict turning points are
\[
0,\quad v,\quad w,\quad N,
\]
and \(c\) is a horizontal inflection point on the increasing middle
branch.  Moreover
\[
H(0)=H(c)=H(N)=0,\qquad
H(v)=-1,\qquad H(w)=1,
\]
and
\[
H''(0)=-\frac2{\alpha^2},\quad
H''(v)=\frac6{\delta^2},\quad
H''(w)=-\frac6{\delta^2},\quad
H''(N)=\frac2{\alpha^2}.
\]
Near \(c\),
\[
H(c+z)=\frac{8}{\delta^3}z^3+O(|z|^4).
\]
\end{proposition}

\begin{proof}
Monotonicity on the five branches is immediate from the definitions.
Near \(t=0\),
\[
I_b(t)=\alpha\sqrt t+O(t),
\]
so inversion gives the same quadratic germ as the outer parabola.
Near \(t=1\),
\[
I_b(t)=b-\delta-\gamma\sqrt{1-t}+O(1-t),
\]
hence inversion gives a quadratic germ with second derivative
\(2/\gamma^2=6/\delta^2\), matching the middle branch at \(v\) and
\(w\).  Finally the middle formula gives
\[
H(c+z)=8z^3/\delta^3+O(|z|^4),
\]
so the first two derivatives match at \(c\) as well.
\end{proof}

\begin{theorem}[Exact real-length chord count]\label{thm:real-count}
For every real \(h>0\),
\[
\nu_h(H_{N,c})
=
2\ind_{\{h<c\}}
+
2\ind_{\{h<e\}}
+
\ind_{\{h=c\}}
+
\ind_{\{h=e\}}
+
\ind_{\{h=N\}}.
\tag{6.2}\label{eq:real-count}
\]
The two indicators at \(c\) and \(e\) are counted separately when
\(c=e\).  Every nonzero-height chord gives a simple zero of the
corresponding difference function.  The zero-height chords are exactly
\[
(0,c),\qquad(c,N),\qquad(0,N),
\]
and they are strict quadratic tangencies of the relevant differences.
\end{theorem}

\begin{proof}
Consider first a negative level \(-t\), \(0<t<1\).  Its three intersection
points, from left to right, are
\[
x_1=-\alpha\sqrt t,\qquad
x_2=I_c(t),\qquad
x_3=c-\delta\theta(t).
\]
Their three pairwise distances are
\begin{align*}
D_1(t)&=I_c(t)+\alpha\sqrt t,\\
D_2(t)&=c-\delta\theta(t)-I_c(t),\\
D_3(t)&=c-\delta\theta(t)+\alpha\sqrt t.
\end{align*}
Put
\[
d_c=c-\delta+\alpha.
\]
Then \(D_1\) is strictly increasing from \(0\) to \(d_c\), while
\(D_2\) is strictly decreasing from \(c\) to \(0\).
For \(D_3\),
\[
\theta'(t)
=
\frac{1}{6t^{2/3}\sqrt{1-t^{1/3}}},
\]
and therefore
\[
2\sqrt t\,\delta\theta'(t)
=
\frac{\delta}
{3\sqrt{t^{1/3}(1-t^{1/3})}}
\ge\frac{2\delta}{3}
>\alpha.
\]
Hence \(D_3\) is strictly decreasing from \(c\) to \(d_c\).

It follows that each \(0<h<c\) occurs exactly twice as a negative-level
distance.  At \(h=d_c\), the endpoint values of \(D_1\) and \(D_3\)
represent the same chord, while \(D_2\) contributes one further chord,
so the count remains two.  At \(h=c\), the only negative-side limit is
the zero-level chord \((0,c)\).  No \(h>c\) occurs on a negative level.

The positive-level picture is the reflection of the same calculation
with \(c\) replaced by \(e\); it contributes two chords for \(0<h<e\),
one zero-level chord at \(h=e\), and none for \(h>e\).  The only
remaining zero-level chord is \((0,N)\).  This proves
\eqref{eq:real-count}.

For a nonzero level, a horizontal separation varying strictly with the
level implies that the two endpoint derivatives are unequal, so the
corresponding zero of \(\Delta_hH\) is simple.  At the merged level
\(t=1\), one endpoint is a turning point and the other has nonzero
derivative, so simplicity still holds.  For the zero-level chords,
both first derivatives vanish, while
\[
(\Delta_cH)''(0)=\frac2{\alpha^2},\qquad
(\Delta_eH)''(c)=\frac2{\alpha^2},\qquad
(\Delta_NH)''(0)=\frac4{\alpha^2}.
\]
Thus all three are strict quadratic tangencies.
\end{proof}

\begin{corollary}[All extremal profiles in \(C^2\)]\label{cor:C2-realize}
For
\[
c=N-r-\frac12,\qquad e=r+\frac12,
\]
the function \(H_{N,c}\) realizes \(V_{N,r}\).
For
\[
c=N-r-1,\qquad e=r+1,
\]
with \(2r+3\le N\), it realizes \(W_{N,r}\).
In the \(V\)-case, every zero below the target is simple.
\end{corollary}

\begin{proof}
Insert the two parameter choices into \eqref{eq:real-count}.  In the
\(V\)-case, neither \(c\) nor \(e\) is an integer.  In the \(W\)-case,
the integer lengths \(e\) and \(c\) each acquire one quadratic tangency,
which changes the adjacent pair from \((2,2)\) to \((3,1)\).
\end{proof}

\section{Polynomialization by a positive derivative weight}

A naive uniform polynomial approximation of \(H_{N,c}\) is insufficient:
it may split a quadratic tangency into two zeros, remove it, or create new
long chords far from the approximation interval.  We instead preserve the
critical-point geometry exactly.

Assume from now on that \(N\in\N\) and that \(c,e=N-c\ge1/2\) are
rational.  Let
\[
J=\left[-\frac14,N+\frac14\right]
\]
and define
\[
\Omega(x)
=
x(x-v)(x-c)^2(x-w)(x-N).
\tag{7.1}\label{eq:Omega}
\]
By Proposition~\ref{prop:Hshape},
\[
U(x):=\frac{H'_{N,c}(x)}{\Omega(x)}
\]
extends continuously across the five zeros of \(\Omega\), and
\[
U(x)>0\qquad(x\in J).
\]
Moreover,
\[
\int_0^c\Omega(x)U(x)\,dx=0,
\qquad
\int_c^N\Omega(x)U(x)\,dx=0.
\tag{7.2}\label{eq:moments}
\]

\begin{lemma}[Positive rational polynomial with two exact moments]
\label{lem:positive-moment}
There exist polynomials \(U_j\in\mathbb{Q}[x]\) such that
\[
U_j(x)>0\quad\text{for all }x\in\R,
\qquad
U_j\longrightarrow U\quad\text{uniformly on }J,
\]
and
\[
\int_0^c\Omega U_j=0,
\qquad
\int_c^N\Omega U_j=0.
\tag{7.3}\label{eq:exact-moments}
\]
\end{lemma}

\begin{proof}
Define the linear map
\[
L(B)=
\left(
\int_0^c\Omega B,\,
\int_c^N\Omega B
\right).
\]
The sign of \(\Omega\) on the four intervals
\[
(0,v),\quad(v,c),\quad(c,w),\quad(w,N)
\]
is, respectively,
\[
-,\quad+,\quad+,\quad-.
\]
Choose four nonnegative continuous bump functions, one supported in each
interval, normalized so that their \(L\)-vectors are
\[
(-1,0),\quad(1,0),\quad(0,1),\quad(0,-1).
\]
Approximate the square root of each bump uniformly on \(J\) by a rational
polynomial \(h_i\).  For sufficiently accurate approximation, the four
vectors
\[
z_i:=L(h_i^2)
\]
remain in the four corresponding quadrants, in cyclic order, with every
adjacent angular gap strictly less than \(\pi\).  Hence the positive cones
generated by adjacent pairs cover \(\R^2\).  Moreover there is a constant
\(C\) such that any \(b\in\R^2\) can be written as
\[
b=\sum_i\lambda_i z_i,\qquad
\lambda_i\ge0,\qquad
\sum_i\lambda_i\le C\|b\|,
\tag{7.4}\label{eq:7.4}
\]
using at most two adjacent generators.  If \(b\in\mathbb Q^2\), the
coefficients \(\lambda_i\) may be chosen rational.

Now approximate \(\sqrt U\) uniformly on \(J\) by
\(p_j\in\mathbb Q[x]\), and choose rational \(\eta_j>0\) with
\(\eta_j\to0\).  Put
\[
B_j=p_j^2+\eta_j.
\]
Then \(B_j>0\) on all of \(\R\), \(B_j\to U\) uniformly on \(J\), and,
by \eqref{eq:moments},
\[
L(B_j)\to0.
\]
Since the endpoints and \(\Omega\) are rational, \(L(B_j)\in\mathbb Q^2\).
Represent \(-L(B_j)\) as in \eqref{eq:7.4} and define
\[
U_j=B_j+\sum_i\lambda_{j,i}h_i^2.
\]
Then \(U_j\in\mathbb Q[x]\), \(U_j>0\) on all of \(\R\), and
\(L(U_j)=0\).  The bound in \eqref{eq:7.4} gives
\(\lambda_{j,i}\to0\), so \(U_j\to U\) uniformly on \(J\).
\end{proof}

\begin{theorem}[Polynomial realization]\label{thm:poly}
Let \(N\in\N\), let \(c,e=N-c\ge1/2\) be rational, and let
\(H_{N,c}\) be the function from \eqref{eq:H}.  Then there exists a
polynomial \(P\in\mathbb Q[x]\) such that
\[
m_n(P)=m_n(H_{N,c})
\qquad\text{for every integer }n\ge1.
\]
Every simple integer-difference zero of \(H_{N,c}\) remains simple for
\(P\), and every integer zero-level tangency remains a double zero.
\end{theorem}

\begin{proof}
Choose \(U_j\) from Lemma~\ref{lem:positive-moment} and define
\[
P_j(x)=\int_0^x\Omega(t)U_j(t)\,dt.
\tag{7.5}\label{eq:Pj}
\]
Then \(P_j\in\mathbb Q[x]\) and the exact moments give
\[
P_j(0)=P_j(c)=P_j(N)=0.
\tag{7.6}\label{eq:7.6}
\]
Also
\[
P_j'(x)=\Omega(x)U_j(x),
\]
so, because \(U_j>0\) on all of \(\R\), the complete real critical set
of \(P_j\) is fixed exactly:
\[
0,\quad v,\quad c,\quad w,\quad N,
\]
with \(c\) a double zero of the derivative and the other four simple
zeros.  Thus no additional real critical point can appear outside the
approximation interval.

Uniform convergence \(U_j\to U\) on \(J\) implies
\[
P_j\to H,\qquad P_j'\to H'
\]
uniformly on \(J\).  Since
\[
H\!\left(-\frac14\right)=-16,\qquad
H(v)=-1,\qquad
H(w)=1,\qquad
H\!\left(N+\frac14\right)=16,
\]
for all sufficiently large \(j\),
\[
P_j\!\left(-\frac14\right)<-2<P_j(v)<0
<P_j(w)<2<P_j\!\left(N+\frac14\right).
\tag{7.7}
\]
The fixed derivative sign pattern now implies that every horizontal level
having at least two intersections with the graph of \(P_j\) has all of its
intersection points in \(J\).  Consequently every horizontal chord of
\(P_j\) has both endpoints in \(J\).  Since
\[
|J|=N+\frac12<N+1,
\]
there are no integer chords of length \(n\ge N+1\).

We next preserve the quadratic tangencies.  The functions \(U_j\) are
uniformly bounded above and away from \(0\) on \(J\).  Since \(\Omega\)
has simple zeros at \(0,N\) and a double zero at \(c\), there exist
constants \(A,B,\varepsilon>0\), independent of sufficiently large \(j\),
such that for \(|t|<\varepsilon\),
\[
-P_j(t)\ge At^2,\qquad
P_j(N+t)\ge At^2,\qquad
|P_j(c+t)|\le B|t|^3.
\]
Shrink \(\varepsilon\) so that \(B\varepsilon<A\).  Then, for
\(0<|t|<\varepsilon\),
\[
\Delta_cP_j(t)>0,\qquad
\Delta_eP_j(c+t)>0,\qquad
\Delta_NP_j(t)>0.
\tag{7.8}
\]
At the central points these differences vanish by \eqref{eq:7.6}.  Their
second derivatives there are strictly positive:
\[
(\Delta_cP_j)''(0)=-P_j''(0)>0,
\]
\[
(\Delta_eP_j)''(c)=P_j''(N)>0,
\]
\[
(\Delta_NP_j)''(0)=P_j''(N)-P_j''(0)>0.
\]
Hence the relevant tangencies remain exactly double and do not split.

It remains to control the other zeros.  For each integer \(1\le n\le N\),
all possible chord endpoints lie in \(J\), so the starting point lies in
the compact interval
\[
J_n=\left[-\frac14,N+\frac14-n\right].
\]
By Theorem~\ref{thm:real-count}, after removing the already treated
tangency neighborhoods, every zero of \(\Delta_nH\) on \(J_n\) is simple
and there are only finitely many.  Choose disjoint neighborhoods of those
zeros on which \((\Delta_nH)'\) has a fixed nonzero sign.  The uniform
\(C^1\)-convergence \(P_j\to H\) gives exactly one simple zero of
\(\Delta_nP_j\) in each such neighborhood.  On the remaining compact set,
\(|\Delta_nH|\) has a positive minimum, so uniform convergence prevents
new zeros.  Since only finitely many integers \(1\le n\le N\) are involved,
one sufficiently large \(j\) works simultaneously for all \(n\).
Take \(P=P_j\).
\end{proof}

\begin{corollary}[Positive-factor form]\label{cor:factor}
The realizing polynomial in Theorem~\ref{thm:poly} may be chosen in the
form
\[
P(x)=x^2(x-N)^2(x-c)^3R(x),
\]
where \(R\in\mathbb Q[x]\) is strictly positive on \(\R\).
\end{corollary}

\begin{proof}
The derivative sign pattern and \eqref{eq:7.6} show that the only real
zeros of \(P\) are \(0,c,N\), with multiplicities \(2,3,2\), respectively.
After dividing out those rational factors, the quotient has no real zero
and is positive at one, hence every, real point.
\end{proof}

\section{Main theorem and transversality}

We now combine necessity and realization.

\begin{theorem}[Extremal classification and rational polynomial realization]
\label{thm:main}
Fix \(N\in\N\).  A nonnegative integer vector
\[
(a_1,\ldots,a_N)
\]
occurs as
\[
(m_1(F),\ldots,m_N(F))
\]
for a continuous \(N\)-extremal function \(F:\R\to\R\) if and only if
it is one of the vectors \(V_{N,r}\) or \(W_{N,r}\) in
\eqref{eq:V}--\eqref{eq:W}.  Every such vector is realizable by a
polynomial in \(\mathbb Q[x]\).
\end{theorem}

\begin{proof}
Necessity is Theorem~\ref{thm:classification}.
For \(V_{N,r}\), choose
\[
c=N-r-\frac12,\qquad e=r+\frac12.
\]
For \(W_{N,r}\), choose
\[
c=N-r-1,\qquad e=r+1.
\]
Corollary~\ref{cor:C2-realize} gives the desired \(C^2\) counts, and
Theorem~\ref{thm:poly} preserves all integer multiplicities in a rational
polynomial.
\end{proof}

\begin{corollary}[Transverse classification]\label{cor:transverse}
The transverse \(N\)-extremal vectors are exactly
\[
V_{N,r},\qquad 0\le r<N/2.
\]
In particular, there are exactly
\[
\left\lceil\frac N2\right\rceil
\]
transverse extremal vectors, and every one has a rational polynomial
realization.
\end{corollary}

\begin{proof}
Every \(V_{N,r}\) is realized by the construction with noninteger
\(c,e\), so all zeros below \(N\) are simple.
Conversely, after the positive tail orientation, every
\(\Delta_nF\) has the same sign at both ends.  If all of its zeros are
simple, their number is even.  Therefore an odd coordinate below \(N\)
is impossible, excluding every \(W\)-vector.
\end{proof}

\section{What the classification changes}

The main theorem replaces three weaker intermediate descriptions.

First, the semigroup theorem records only which lengths may be absent.
At extremality, Theorem~\ref{thm:terminal} collapses this arithmetic
freedom to one integer \(r\): every missing support is a terminal block.

Second, terminal support alone permits formal central allocations
\[
(1,3),\quad(2,2),\quad(3,1)
\]
in each complementary pair.  Theorem~\ref{thm:classification} shows that
these choices are not independent.  Global compatibility among the
difference functions allows at most one odd pair, and it must sit
immediately next to the \((4,0)\) terminal band.

Third, a finite list of explicit low-degree examples cannot settle
existence for all \(N\).  The inverse-branch formula
\eqref{eq:real-count} solves the realization problem before any polynomial
is chosen: it prescribes the full continuum of chord lengths.  The
polynomialization theorem then transfers the exact integer data while
locking the critical-point geometry on the entire real line.  Thus no
finite numerical verification is needed in the final classification.

\section{Further problems}

The extremal distinct-chord problem is now closed, but several natural
questions remain.

\begin{problem}[Minimum degree]
Determine the least degree of a polynomial realizing each
\(V_{N,r}\) and \(W_{N,r}\).
\end{problem}

\begin{problem}[Uniform degree bound]
Does there exist an absolute integer \(D\) such that every extremal vector,
for every target \(N\), has a polynomial realization of degree at most \(D\)?
The existence proof in Theorem~\ref{thm:poly} gives finite degree for each
parameter pair but no uniform bound.
\end{problem}

\begin{problem}[Realization spaces]
For a fixed \(N\) and a fixed \(V\)- or \(W\)-profile, describe the topology
and geometry of the realization space modulo translation and positive
vertical scaling.  In particular, understand how a \(W\)-profile sits on
the discriminant boundary between transverse \(V\)-type regions.
\end{problem}

\begin{problem}[Beyond the minimum]
Suppose \(m_N>0\) but \(M(F)>2N-1\).  Determine sharp propagation laws
between missing lengths and excess multiplicities, and classify the first
few levels above the minimum.
\end{problem}

\section{Conclusion}

For integer horizontal chords, the sharp total \(2N-1\) is only the first
layer of the problem.  Equality exposes a rigid interaction between the
additive structure of missing lengths and the zero structure of the
difference functions \(\Delta_nF\).

The complete picture is now simple.  Missing lengths form one terminal
interval.  Every central complementary pair is balanced as \((2,2)\),
except that one adjacent pair may degenerate to \((3,1)\).  These are
exactly the \(V\)- and \(W\)-families, and no other extremal vector is
possible.  Conversely, every one of these vectors is realized by a
rational polynomial.

The main constructive idea is to work in inverse-branch coordinates:
horizontal chord multiplicity is then controlled by monotone separation
functions rather than by solving a family of high-degree difference
equations.  The subsequent factorization
\[
P'=\Omega U,\qquad U>0,
\]
with two exact moment conditions, preserves the branch geometry while
moving into the algebraic category.  This combination of global
difference constraints, inverse-branch geometry, and positive-weight
polynomialization is the structural core of the classification.

\begin{thebibliography}{99}

\bibitem{Levy1934}
P.~L\'evy,
\newblock Sur une g\'en\'eralisation du th\'eor\`eme de Rolle,
\newblock \emph{Comptes Rendus de l'Acad\'emie des Sciences, Paris}
\textbf{198} (1934), 424--425.

\bibitem{Hopf1937}
H.~Hopf,
\newblock \"Uber die Sehnen ebener Kontinuen und die Schleifen
geschlossener Wege,
\newblock \emph{Commentarii Mathematici Helvetici} \textbf{9}
(1936/37), 303--319.

\bibitem{Oxtoby1972}
J.~C. Oxtoby,
\newblock Horizontal chord theorems,
\newblock \emph{American Mathematical Monthly} \textbf{79} (1972),
468--475.
\newblock \url{https://doi.org/10.1080/00029890.1972.11993068}.

\bibitem{BurnsDavidovichDavis2017}
K.~Burns, O.~Davidovich, and D.~Davis,
\newblock Average pace and horizontal chords,
\newblock \emph{Mathematical Intelligencer} \textbf{39}, no.~4
(2017), 41--45.
\newblock \url{https://doi.org/10.1007/s00283-017-9713-2}.

\bibitem{DavisTroubetzkoy2024}
D.~Davis and S.~Troubetzkoy,
\newblock The horizontal chord set,
\newblock \emph{American Mathematical Monthly} \textbf{131}, no.~6
(2024), 519--525.
\newblock \url{https://doi.org/10.1080/00029890.2024.2325862}.

\bibitem{CiudinIonascu2026}
I.~Ciudin and E.~J. Ionascu,
\newblock Universal Chord Theorem and a Topological Analysis,
\newblock arXiv:2601.14120, 2026.
\newblock \url{https://arxiv.org/abs/2601.14120}.

\bibitem{VSO2025}
Central Subject-Methodological Commission in Mathematics,
\newblock Materials for the final stage of the 51st All-Russian
Mathematical Olympiad, second day, Problem 11.8 and official solution,
\newblock Sirius, April 16--22, 2025.
\newblock Problem proposed by A.\,S.~Kuznetsov.
\newblock
\url{https://olympiads.mccme.ru/vmo/2025/final/sol2.pdf}.

\end{thebibliography}

\end{document}
